\section{Properties at the checked revision}
\label{sec:properties}
Each reference below uses the full pin in Appendix~\ref{sec:artifact} and its
verbatim axiom line. Unless stated otherwise, economic results quantify over
an arbitrary reachable world and an arbitrary transaction accepted by the
relevant judgement, not just a builder's output. The distinction between the
original registry fold and the economic observation is explicit.

\paragraph{1. Fixed parameters.}
\checked An accepted original registry step or batch satisfies
\decl{Singular.onlyRootChanged}: the root may change but all other configuration
fields, including $R,t_{\min},\delta,W_p,W_r$ and policy bindings, are preserved.
\evidence{parameters_preserved_by_fold}
\evidence{parameters_preserved_by_batch}
\hand The economic fold reuses that transition through
\decl{Singular.foldProjectedRegistry} and \decl{Singular.foldBatch}. Inspection of
\decl{Singular.foldAccepted} shows that its resulting registry comes from this
batch. The projection restores carried termination holdings for the original
registry interface and supplies synthetic approval evidence to that interface;
it is not an actual termination application tag or a second mint.
There is no registry update or migration operation in the defined transition
set. This is a statement about that set, not an assertion that an arbitrary
deployed policy cannot spend the state.

\paragraph{2. Rejection preserves the registry.}
\checked Accepted rejection leaves the registry, pending terminations and
singular supply equal. Its observed spending and outputs carry no state or
singular tokens; no such mint is present. Request-tag burns are separate.
The genuine referenced state remains unspent after applying the transaction.
\evidence{rejection_leaves_registry_untouched}
\evidence{rejection_requires_referenced_state}

\paragraph{3. State noninterference, with its actual scope.}
\checked If the \emph{original admitted registry fold} was valid before a
rejection, it remains valid afterward. A supplied list of needed unspent
outputs also remains present if all its other inputs are disjoint from the
rejection's spends. The referenced state is protected by the rejection checks.
\evidence{rejection_cannot_invalidate_fold}
\missing This declaration preserves \decl{Singular.admittedFoldBatch} and
needed outputs. It is not a theorem preserving the whole later
\decl{Singular.foldAccepted} judgement, fees, wallet state or ledger acceptance.
The shared-input counterexample below prevents dropping disjointness.

\paragraph{4. Live protection and a non-backdated clock.}
\checked A pending registration whose key still admits registration cannot
be a member of an accepted rejection whose interval begins before $s+W_p+W_r$.
This includes the reclaim window, not just the processing window. Given an
actual inclusion time $\tau\in[L,U)$, the checked booking rule gives
$\tau\le s$ and $\tau+W_p+W_r\le s+W_p+W_r$.
\evidence{live_request_not_rejected}
\evidence{booked_time_not_before_inclusion}
A current mismatch is outside the compatibility hypothesis: a losing race
can be rejected immediately. The inequality states an expiry threshold, not
an inclusion guarantee or the existence of a retraction transaction.

\paragraph{5. Current evidence is mandatory.}
\checked An accepted rejection has an unspent reference with the exact state
identity, inline configuration and current root; that reference is not spent.
\evidence{rejection_requires_referenced_state}
The per-member mismatch test reads this current state. An alleged booking-root
fault is neither required nor substituted. This rules out evidence-free
acceptance in the economic judgement; it does not verify concrete Merkle code.

\paragraph{6. Refunds are additive per address.}
\checked Rejected members and their input references are distinct. For every
refund address $A$, outputs paid to it total at least
\[
 F_A=\sum_{q:\,\mathrm{refund}(q)=A}(H_q-t_q).
\]
The obligations have one entry per address. The output cannot be counted as a
refund at a different address, and duplicated members cannot reuse it.
\evidence{rejection_batch_members_distinct}
\evidence{rejection_refund_floor}
\evidence{refunds_sum_per_address}
\checked Accepted retraction identifies one pending registration, its signed
refund address and a contained reclaim interval, burns both tags, returns at
least its entire $H$, reads and spends no state output, and removes that
request from the pool.
\evidence{retraction_refund_floor}
\checked The request is among the spending inputs; ordinary funding may
accompany it. The retraction balances, pays its fee from that funding, donates
zero and leaves the treasury unchanged. Value at other addresses is bounded
by ordinary funding minus the fee. Neither the cost nor the tip can pay it.
\evidence{retraction_pays_only_the_refund_and_the_fee}

\paragraph{7. Charge and return.}
\checked An economic fold registering $n$ keys donates at least $nR$;
Equation~\eqref{eq:donation} gives its exact amount. The charge is included in
the rejected address's refund floor and in the retracted request's floor.
\evidence{fold_donates_registration_cost}
\evidence{registration_cost_returned_on_exit}
\hand A successful fold removes its registration request, and the lifecycle
has no operation returning its treasury charge. This concerns the mechanism's
own transfers, not a prohibition on a later treasury payment to the same party.

\paragraph{8. Terminations are unrejectable.}
\checked Every genuine pending termination keeps its key active. Mismatch
therefore fails and there is no expiry ground. It cannot occur as a member of
any accepted rejection, whichever ground is offered.
\evidence{termination_unrejectable}

\paragraph{9. At most one pending termination per key.}
\checked Singular supply counts application custody \emph{plus} actual tags
in pending terminations. It is one exactly for an active key and zero otherwise,
never greater than one. Carried count equals pending termination count, so at
most one is pending for each key.
\evidence{world_supply_sync}
\evidence{at_most_one_pending_termination}
Counting only application custody after release would give the wrong invariant.

\paragraph{10. Termination provenance.}
\checked Every termination selected by an accepted economic fold was pending,
is observed in the unspent set, carries its key's genuine singular tag, and
has an earlier accepted termination booking from a reachable world. That
booking consumed the unique matching custody holding. The application's
release rule itself is an external premise, not rechecked by Singular.
\evidence{only_booked_terminations_fold}
\evidence{termination_request_carries_genuine_token}

\paragraph{11. Conditional model availability.}
\checked For any pending termination in a reachable world, any natural-number
time $u$, funding $I$ and fee $f$ with $f\le I$ and $f\le t$, and a predicate
saying a folder accepts its tip, there exist a coherent unspent set and an
accepted singleton fold with interval $[u,u+1)$. It ends the key as Terminal,
removes the request and produces a reachable world.
\evidence{booked_termination_foldable}
\checked For a pending registration in a reachable world, choose a nonempty
interval wholly inside its booked reclaim window, signers including its
refund address, and ordinary funding $I\ge f$. There exist a coherent unspent
set, a balanced accepted retraction with exactly that interval, signers, fee
and funding total, and a reachable after-world without the request. The full
held lovelace is returned. The theorem permits a positive fee and does not
require it to fit inside the refundable tip.
\evidence{pending_registration_can_be_retracted}
\missing The theorem constructs its own unspent set; it does not find funding
or prove an accepted transaction over an arbitrary live ledger snapshot.
Its witness includes a zero-lovelace state output, illustrating the absent
ledger-minimum premise. It proves no willingness, inclusion or bounded delay.

\paragraph{12. Processor cap is a fold allocation bound.}
\checked For an accepted economic fold, $O\le I+T$, with
$O+f=I+T$ and $f\le I$ as above.
\evidence{processors_take_at_most_the_tips}
\hand Thus net ordinary increase is at most $T$. The formal theorem is about
ordinary outputs and inputs, not ownership, every address receiving funds or
every exit kind. With zero ordinary funding the fee condition forces $f=0$.
No exact rejection reward, positive wage or service incentive follows.

\paragraph{13. A measured minimum tip.}
\proposed Choose $t_{\min}$ to cover the maximum measured single-request exit
cost in a declared envelope, including the reference read and required refund
output. The derivation and assumptions follow in Section~\ref{sec:sizing}.
\missing No checked declaration or completed measurement establishes this
cost cover. Retraction's modeled fee funding is explicit, but a measured
ledger envelope and implementation correspondence remain open. Its refunded
tip is not a processor payment and does not enter a processor-income floor.

\paragraph{Supporting supply statement.}
\checked In every reachable world, the number of trie-ready tags actually
held in the two pending pools equals the total number of pending requests.
Application tags held in pending registrations equal their number, with
exactly one in each. This is not a total application-token supply theorem.
\evidence{tag_supplies_equal_pending}
The world has no separate external trie-ready circulation; lifting its counted
supply to all ledger assets requires mint/spend correspondence. Arbitrary
application minting elsewhere is explicitly outside the counted domain.
